%! Multiplying \& Dividing Rational Expressions — Unit 4, Lesson 4.2 — Slides (Beamer)
\documentclass[aspectratio=169, 11pt]{beamer}
\usepackage{algebra2-beamer}   % provides \forestheader{} and \sectionlabel[color]{}
% Standard coordinate grid + axes: {xmin}{xmax}{ymin}{ymax}
\newcommand{\lgrid}[4]{%
  \draw[step=1, linegray!45, thin] (#1,#3) grid (#2,#4);
  \draw[->, thick, charcoal] (#1,0)--(#2,0) node[below right, font=\tiny]{$x$};
  \draw[->, thick, charcoal] (0,#3)--(0,#4) node[left, font=\tiny]{$y$};}

\begin{document}

% TITLE SLIDE (CourseName is not defined in beamer, so the name is literal)
{
\setbeamercolor{background canvas}{bg=forest}
\begin{frame}[plain]
  \vspace{2.8cm}\hspace{0.55cm}%
  \begin{minipage}{0.9\textwidth}
    {\color{goldacc}\bfseries\small LESSON 4.2}\\[0.5cm]
    {\color{white}\bfseries\LARGE Multiplying \& Dividing Rational Expressions}\\[1.2cm]
    {\color{forestmid}\small Algebra 2: Shepherd~$\cdot$~Unit 4: Rational Functions}
  \end{minipage}
\end{frame}
}

% HOOK
\begin{frame}[t, plain]
  \forestheader{Same expression. One disagreement.}
  \sectionlabel{Hook}
  \begin{columns}[T]
    \begin{column}{0.42\textwidth}
      \[ Q(x)=\frac{x+1}{x-5}\div\frac{x-2}{x+3} \]
      \[ R(x)=\frac{(x+1)(x+3)}{(x-5)(x-2)} \]
      \vspace{-8pt}
      \begin{itemize}\setlength{\itemsep}{4pt}
        \item $R$ is $Q$ after keep--change--flip.
      \end{itemize}
    \end{column}
    \begin{column}{0.56\textwidth}
      \centering
      \renewcommand{\arraystretch}{1.4}
      \begin{tabular}{|l|c|c|c|c|}
      \hline
      $x$ & $0$ & $1$ & $2$ & $-3$ \\ \hline
      $Q(x)$ & $\frac{3}{10}$ & $2$ & \textbf{undef.} & \textbf{undef.} \\ \hline
      $R(x)$ & $\frac{3}{10}$ & $2$ & \textbf{undef.} & $0$ \\ \hline
      \end{tabular}
    \end{column}
  \end{columns}
  \vspace{6pt}
  \begin{block}{Today}
    Yesterday, \textbf{cancelling} hid a restriction. Today, \textbf{flipping} hides a different one.
  \end{block}
\end{frame}

% MULTIPLYING
\begin{frame}[t, plain]
  \forestheader{Multiplying: the rule is easy, the discipline is factoring}
  \sectionlabel{Products}
  \begin{columns}[T]
    \begin{column}{0.50\textwidth}
      \[ \frac{a}{b}\cdot\frac{c}{d}=\frac{a\,c}{b\,d}, \qquad b\neq0,\ d\neq0 \]
      \begin{itemize}\setlength{\itemsep}{4pt}
        \item $\dfrac{6}{35}\cdot\dfrac{25}{9}$: factor, divide out, \emph{then} multiply
              $\Rightarrow \dfrac{10}{21}$.
        \item Multiply first and you get $\dfrac{150}{315}$ --- and still have to reduce it.
      \end{itemize}
    \end{column}
    \begin{column}{0.46\textwidth}
      \begin{block}{Why it matters more with letters}
        An expanded product \textbf{hides its own factors}. Unit 3 showed you what it costs to get
        them back.
      \end{block}
      \vspace{3pt}
      {\footnotesize \textbf{Factor everything first. Divide out. Multiply last --- and often not at
      all.}}
    \end{column}
  \end{columns}
\end{frame}

% FOUR STEPS + ANCHOR
\begin{frame}[t, plain]
  \forestheader{Four steps --- now with two denominators}
  \sectionlabel[goldacc]{Procedure}
  \begin{columns}[T]
    \begin{column}{0.48\textwidth}
      \begin{enumerate}\setlength{\itemsep}{5pt}
        \item \textbf{Factor} every numerator and denominator.
        \item \textbf{List the restrictions} --- from \emph{both} original denominators.
        \item \textbf{Divide out} across the whole product.
        \item \textbf{Write} it, restrictions attached, still factored.
      \end{enumerate}
    \end{column}
    \begin{column}{0.50\textwidth}
      \begin{block}{The anchor}
        \[ \frac{x^2-9}{x^2+2x-8}\cdot\frac{x+4}{x^2+3x} \]
        \[ =\frac{(x-3)(x+3)}{(x+4)(x-2)}\cdot\frac{x+4}{x(x+3)}
           =\frac{x-3}{x(x-2)} \]
        \centering $x\neq 0,\;2,\;-3,\;-4$
      \end{block}
      \vspace{2pt}
      {\footnotesize Two of the four restrictions are invisible in the answer.}
    \end{column}
  \end{columns}
\end{frame}

% DIVIDING
\begin{frame}[t, plain]
  \forestheader{Dividing: keep, change, flip --- \emph{then} cancel}
  \sectionlabel{Quotients}
  \begin{columns}[T]
    \begin{column}{0.48\textwidth}
      \[ \frac{a}{b}\div\frac{c}{d}=\frac{a}{b}\cdot\frac{d}{c} \]
      \begin{itemize}\setlength{\itemsep}{4pt}
        \item Dividing by something $=$ multiplying by what undoes it, and
              $\frac{c}{d}\cdot\frac{d}{c}=1$.
        \item A fraction has a reciprocal \textbf{only when its numerator is nonzero}.
      \end{itemize}
    \end{column}
    \begin{column}{0.48\textwidth}
      \begin{block}{Cancelling across a $\div$ is illegal}
        Wrong: $\dfrac{4}{3}\div\dfrac{5}{4}\to\dfrac{1}{3}\div\dfrac{5}{1}=\dfrac{1}{15}$ \\[3pt]
        Right: $\dfrac{4}{3}\cdot\dfrac{4}{5}=\dfrac{16}{15}$
      \end{block}
      \vspace{3pt}
      {\footnotesize The $\div$ symbol is not a fraction bar. Until you flip, nothing is lined up.}
    \end{column}
  \end{columns}
\end{frame}

% THE BIG IDEA
\begin{frame}[t, plain]
  \forestheader{Three sources of a restriction}
  \sectionlabel{The big idea}
  \begin{center}
    \[ \frac{A}{B}\div\frac{C}{D} \]
  \end{center}
  \vspace{-4pt}
  \begin{columns}[T]
    \begin{column}{0.32\textwidth}
      \begin{block}{1. $B\neq0$}
        The first fraction must exist. \\[3pt]
        {\footnotesize Visible before \emph{and} after.}
      \end{block}
    \end{column}
    \begin{column}{0.32\textwidth}
      \begin{block}{2. $D\neq0$}
        The divisor must exist. \\[3pt]
        {\footnotesize Invisible \emph{after} the flip --- $D$ moves upstairs.}
      \end{block}
    \end{column}
    \begin{column}{0.32\textwidth}
      \begin{block}{3. $C\neq0$}
        You cannot divide by zero. \\[3pt]
        {\footnotesize Invisible \emph{before} the flip --- $C$ is upstairs already.}
      \end{block}
    \end{column}
  \end{columns}
  \vspace{6pt}
  \begin{center}
    \textbf{Flipping trades one blind spot for another.} Read the restrictions off the problem as
    written --- then add the divisor's numerator.
  \end{center}
\end{frame}

% DIVISION ANCHOR
\begin{frame}[t, plain]
  \forestheader{The division anchor, all three sources}
  \sectionlabel[goldacc]{Worked}
  \begin{columns}[T]
    \begin{column}{0.54\textwidth}
      \[ \frac{x^2-4}{x^2+7x+12}\div\frac{x^2+2x}{x+3} \]
      \[ =\frac{(x-2)(x+2)}{(x+3)(x+4)}\cdot\frac{x+3}{x(x+2)}
         =\frac{x-2}{x(x+4)} \]
    \end{column}
    \begin{column}{0.44\textwidth}
      \begin{block}{Where each one came from}
        $x\neq-3,-4$ \; first denominator \\[2pt]
        $x\neq-3$ \; divisor's denominator \\[2pt]
        $x\neq0,-2$ \; divisor's \textbf{numerator}
      \end{block}
      \vspace{3pt}
      {\footnotesize Final: $x\neq0,-2,-3,-4$. Only two of them are visible in the answer.}
    \end{column}
  \end{columns}
\end{frame}

% THE PICTURE
\begin{frame}[t, plain]
  \forestheader{What the restrictions look like}
  \sectionlabel{The picture}
  \begin{columns}[T]
    \begin{column}{0.44\textwidth}
      \centering
      \begin{tikzpicture}[scale=0.34]
        \lgrid{-4}{5}{-5}{4}
        \begin{scope}
          \clip (-4,-5) rectangle (5,4);
          \draw[very thick, forest] (-4,-5)--(5,4);
        \end{scope}
        \draw[forest, very thick, fill=white] (0,-1) circle (5pt);
        \draw[forest, very thick, fill=white] (-1,-2) circle (5pt);
      \end{tikzpicture}
    \end{column}
    \begin{column}{0.54\textwidth}
      \[ P(x)=\frac{x^2-1}{x}\cdot\frac{x}{x+1}=x-1 \]
      \begin{block}{One hole per denominator}
        $x\neq0$ (first fraction) $\Rightarrow$ hole at $(0,-1)$ \\[3pt]
        $x\neq-1$ (second fraction) $\Rightarrow$ hole at $(-1,-2)$
      \end{block}
      \vspace{3pt}
      {\footnotesize The answer $x-1$ has no denominator left to show them --- and they are still
      excluded.}
    \end{column}
  \end{columns}
\end{frame}

% ERROR GALLERY
\begin{frame}[t, plain]
  \forestheader{The four ways to lose a point today}
  \sectionlabel{Watch out}
  \begin{columns}[T]
    \begin{column}{0.48\textwidth}
      \begin{block}{Multiplying out first}
        You get a quartic nobody wants to re-factor. Factor \emph{first}.
      \end{block}
      \vspace{3pt}
      \begin{block}{Cancelling before flipping}
        $\dfrac{x}{4}\div\dfrac{5}{x}=\dfrac{x^2}{20}$, not $\dfrac{1}{20}$.
      \end{block}
    \end{column}
    \begin{column}{0.48\textwidth}
      \begin{block}{Flipping the wrong fraction}
        Keep the \emph{first}. Flip the \emph{divisor}.
      \end{block}
      \vspace{3pt}
      \begin{block}{Missing the divisor's numerator}
        Today's signature error --- and the restriction nothing in the answer will remind you of.
      \end{block}
    \end{column}
  \end{columns}
\end{frame}

% ROADMAP
\begin{frame}[t, plain]
  \forestheader{Where this goes}
  \sectionlabel{Connections}
  \textbf{Today:} factor, flip, divide out --- and know all three places a restriction can come
  from.\\[8pt]
  \begin{itemize}\setlength{\itemsep}{3pt}
    \item \textbf{4.3} Adding \& subtracting --- and a \emph{complex fraction}, which is just
          combine-then-divide, i.e.\ today.
    \item \textbf{4.5} Graphing from the equation: every restriction you list is a hole or a wall.
    \item \textbf{4.6} Solving rational equations --- an extraneous solution is a restriction that
          sneaked back in.
  \end{itemize}
  \vspace{6pt}
  \begin{center}
    \textbf{Flipping trades one blind spot for another.} Say it once more on the way out.
  \end{center}
\end{frame}

\end{document}
